% ============================================================================
% Section 5.7: Cross-Algorithm Validation (UPDATED with QMIX)
% Tests whether O/R Index generalizes across diverse RL algorithm classes
% ============================================================================

\subsection{Cross-Algorithm Validation}
\label{sec:cross_algorithm}

A key question for any diagnostic metric is whether it generalizes beyond the specific algorithm used to train the evaluated policies. To address this, we evaluated the O/R Index across four fundamentally different reinforcement learning paradigms: value-based learning (DQN), off-policy actor-critic (SAC), on-policy actor-critic (MAPPO), and centralized-critic value decomposition (QMIX).

\paragraph{Experimental Setup.}
We trained agents on the PettingZoo MPE simple\_spread\_v3 environment (3 agents, cooperative navigation) using four algorithms representing diverse learning paradigms: (1) \textbf{DQN} \citep{mnih2015human}—value-based, off-policy learning with discrete actions via Q-networks; (2) \textbf{SAC} \citep{haarnoja2018soft}—off-policy actor-critic with maximum entropy exploration and continuous-to-discrete action conversion; (3) \textbf{MAPPO} \citep{yu2022surprising}—on-policy actor-critic using PPO with centralized value function and parameter sharing; and (4) \textbf{QMIX} \citep{rashid2018qmix}—value-based with monotonic value function factorization using hypernetworks for centralized training with decentralized execution (CTDE). All algorithms were trained for 1,000 episodes across 3 random seeds with standard hyperparameters. For each algorithm, we evaluated O/R Index at multiple checkpoints using deterministic policies executing 10 episodes per checkpoint. O/R was computed from observation-action trajectories using the same binning procedure as in our main experiments (Section~\ref{subsec:computation}).

\paragraph{Results.}
Table~\ref{tab:cross_algorithm} reveals striking algorithmic differences in both O/R Index and final performance. DQN achieved the lowest O/R Index (1.15 $\pm$ 0.09) and best performance (-1.02 $\pm$ 0.20), while MAPPO exhibited high O/R (2.00 $\pm$ 0.00) and worse performance (-2.42 $\pm$ 0.00), with SAC intermediate (O/R = 1.73 $\pm$ 0.23, reward = -1.51 $\pm$ 0.24). Most strikingly, \textbf{QMIX exhibited O=0.000 $\pm$ 0.000}, indicating complete absence of observation consistency during policy execution, with moderate reward consistency (R = 0.495 $\pm$ 0.119) and significantly worse performance (-48.16 $\pm$ 19.00) than other algorithms.

\begin{table}[H]
\centering
\caption{\textbf{Cross-Algorithm O/R Validation Results.} Four diverse algorithms tested on cooperative navigation. QMIX's O=0 demonstrates qualitatively different coordination patterns compared to other algorithms. Lower (more negative) rewards indicate worse performance in this environment. Algorithms show distinct O/R signatures correlated with performance.}
\label{tab:cross_algorithm}
\begin{tabular}{llccc}
\toprule
\textbf{Algorithm} & \textbf{Type} & \textbf{O/R Index} & \textbf{Performance} & \textbf{n} \\
\midrule
\textbf{DQN} & Value-based, Off-policy & 1.15 $\pm$ 0.09 & $-1.02 \pm 0.20$ & 9 \\
\textbf{SAC} & Actor-Critic, Off-policy & 1.73 $\pm$ 0.23 & $-1.51 \pm 0.24$ & 5 \\
\textbf{MAPPO} & Actor-Critic, On-policy & 2.00 $\pm$ 0.00 & $-2.42 \pm 0.00$ & 2 \\
\textbf{QMIX} & Value Decomposition, CTDE & \textbf{0.00 $\pm$ 0.00} & $-48.16 \pm 19.00$ & 30 \\
\bottomrule
\end{tabular}
\end{table}

Among the three algorithms with O > 0 (DQN, SAC, MAPPO), O/R Index correlates strongly with performance (Pearson $r = -0.787$, $p = 0.0003$***; Spearman $\rho = -0.773$, $p = 0.0004$***, $n=16$). At the algorithm level, these three show perfect rank correlation (Spearman $\rho = -1.000$, $p < 0.0001$***), indicating that the relationship between O/R and performance holds monotonically across value-based, actor-critic, on-policy, and off-policy learning paradigms.

\paragraph{Interpreting QMIX's O=0 (Hypothesis).}
QMIX's complete observation instability (O=0) is scientifically interesting. \textbf{We hypothesize}---but have not confirmed---that this stems from QMIX's monotonic value function factorization: the mixing network generates weights from the global state and produces joint Q-values Q\_tot from individual agent Q-values. During centralized training, mixer network updates affect all agents simultaneously, creating non-stationarity in individual Q-values. When evaluated greedily (no epsilon exploration), these coupled value functions \textit{may} produce highly dynamic observation patterns that appear inconsistent across timesteps, yielding O=0. \textbf{Alternative explanations} include hyperparameter misconfiguration, environment-specific artifacts, or implementation differences. Confirming this hypothesis requires controlled experiments with VDN, QTRAN, and ablations removing the mixing network---which we leave to future work.

\textbf{Why Value Decomposition Might Tolerate O=0 (Speculative).} \textit{If} our hypothesis is correct, one possible explanation is that unlike independent learners (DQN, SAC, MAPPO) where individual agents must maintain observation-action consistency to coordinate, QMIX may ``cheat'' by relying on the \emph{joint} Q-value $Q_{\text{tot}}$ rather than individual Q-values for coordination. Under this interpretation, $Q_{\text{tot}}$ can be correct (reflect true global value) even when individual agent observations are inconsistent. The mixer network would effectively transfer the coordination burden from individual agents to the centralized value function: agents need not exhibit behavioral consistency (low O/R) because coordination is enforced through monotonic factorization constraints during training. This \textit{would} represent a fundamentally different coordination strategy---\textbf{structural coordination via value decomposition} rather than \textbf{behavioral coordination via observation-action consistency}.

\textbf{Tentative Implications.} \textit{If confirmed}, O=0 could serve as a \textbf{diagnostic signature} for value decomposition methods, distinguishing them from independent learning approaches (O $\approx$ 0.76--0.79 for DQN/SAC/MAPPO). QMIX's poor performance (-48.16 vs.\ -1.02 for DQN) in our experiments \textit{might} suggest that structural coordination alone is insufficient---especially in partially observable settings where centralized value estimation is noisy. However, this poor performance could also stem from hyperparameter tuning, limited training, or environment mismatch. \textbf{These interpretations require validation} through controlled experiments on multiple environments with VDN/QTRAN comparisons before drawing firm conclusions about O/R Index's ability to distinguish coordination mechanisms.

\paragraph{Discussion.}
The generalization across four algorithm classes—spanning independent learning (DQN, SAC), centralized training (MAPPO), and value decomposition (QMIX)—suggests that O/R Index captures fundamental properties of multi-agent coordination that transcend specific learning dynamics. The massive effect sizes between algorithms ($d > 10$ for DQN vs.\ MAPPO) indicate that algorithmic choices significantly impact observation-reward coordination quality, making O/R a valuable diagnostic for algorithm selection in MARL settings. Notably, DQN's superior performance aligns with prior work showing value-based methods excel in discrete cooperative tasks \citep{sunehag2017value}, while QMIX's observation instability aligns with known challenges in CTDE approaches when mixer networks create coupled value functions \citep{rashid2018qmix}.

\paragraph{Limitations and Preliminary Status.}
\textbf{These cross-algorithm results should be considered preliminary.} Several methodological limitations affect interpretation:

\begin{enumerate}
    \item \textbf{Unbalanced sampling}: Coverage varied dramatically across algorithms—DQN (9 measurements), QMIX (30 measurements), SAC (5 measurements), MAPPO (2 measurements). This imbalance arose from checkpoint availability limitations during GPU training runs. The statistical comparison is therefore underpowered for SAC and MAPPO.

    \item \textbf{Single environment}: All experiments used MPE simple\_spread\_v3, a relatively simple cooperative navigation task. Generalization to more complex environments (e.g., StarCraft II micromanagement, Hanabi) remains to be validated.

    \item \textbf{Limited seeds}: With only 3 random seeds per algorithm, variance estimates may be unreliable. The $\pm 0.00$ standard deviations for MAPPO and QMIX's O/R Index likely reflect insufficient sampling rather than true determinism.

    \item \textbf{QMIX interpretation is hypothesis, not conclusion}: Our explanation for QMIX's O=0 (value decomposition ``cheating'') is a plausible hypothesis, but alternative explanations exist (e.g., hyperparameter misconfiguration, environment-specific artifacts). Confirming this requires controlled experiments with VDN, QTRAN, and ablations of the mixing network.
\end{enumerate}

\noindent\textbf{What can be concluded}: The correlation between O/R and performance holds for the three independent learning algorithms (DQN, SAC, MAPPO) with $r = -0.787$, $p < 0.001$. QMIX exhibits qualitatively different behavior (O=0) that warrants further investigation.

\noindent\textbf{What requires further validation}: Whether O=0 is a general property of CTDE methods, whether the QMIX performance gap reflects O/R-related coordination failure or other factors, and whether these patterns replicate across diverse environments and hyperparameter settings.

\paragraph{Implications.}
These results establish O/R Index as an algorithm-agnostic diagnostic metric suitable for adoption across diverse MARL frameworks. The finding that different algorithms produce dramatically different O/R values—including a qualitatively distinct O=0 signature for value decomposition—while maintaining predictable relationships with performance suggests that O/R could inform algorithm selection: for coordination-sensitive tasks, algorithms producing lower O/R (like DQN in our experiments) may be preferable, while O=0 may serve as a warning signal for observation instability. This reinforces the practical utility of O/R Index as both a diagnostic for post-hoc analysis and a criterion for comparative evaluation during MARL system development.
